The rung summand, its mode vectors and the involution $*$ are those of Fact \claimref{standard-model.imported-tower-commuting-copies-fock-module}: the adjoined null pairs satisfy $v_{a}w_{b} + w_{b}v_{a} = \delta_{ab}$, and $*$ fixes the positive frame generators and negates the negative ones, so on a vector $x^{*} = \theta(\bar{x})$. In the basis $(v_{4}, v_{5}, w_{4}, w_{5})$ the metric is $\ip{v_{a}}{w_{b}} = \tfrac12\delta_{ab}$, all other pairings zero, and $\theta$ exchanges $v_{a}$ and $w_{a}$.

(1) From the anticommutators, $[v_{a}w_{b}, v_{c}] = \delta_{bc}v_{a}$ and $[v_{a}w_{b}, w_{c}] = -\delta_{ac}w_{b}$: $E_{ab}$ is the matrix unit on the $v$'s and minus its transpose on the $w$'s. These are real operators on $\mathcal{R}$, each preserving $Q$ (each $E_{ab}$ with $a \neq b$ is the real bivector $v_{a}\wedge w_{b}$, and $E_{aa} - \tfrac12$ is $v_{a}\wedge w_{a}$), so their real span is $\gl(2,\R)$ inside $\so(2,2)$. A real combination $\sum c_{k}G_{k}$ of the four $\uu(2)$ generators preserves $\mathcal{R}$ if and only if its matrix is real; the imaginary parts come only from the three generators carrying $I_{4}$ and are independent, so the coefficients of those three vanish and the solution space is spanned by $E_{45} - E_{54}$. The rung number $N_{R} = E_{44} + E_{55}$ acts as $+1$ on the $v$'s and $-1$ on the $w$'s, so the hypercharge $3N_{R} - 2N_{W}$ of Proposition \claimref{standard-model.ten-transported-ambient-space-doublets-are} has the real weights $\pm 3$ on rung vectors, and its compact $U(1)$, $\exp(I_{4}\varphi\,y)$, multiplies $v_{a}$ by $e^{3I_{4}\varphi}$, which leaves $\mathcal{R}$ only for $\varphi \in \tfrac{\pi}{3}\mathbb{Z}$.

(2) The six bivectors $e_{i}e_{j}$ span $\so(2,2)$ acting by commutator; the two whose generators have equal square are rotations with $\tr(X^{2}) < 0$, the four mixed ones boosts with $\tr(X^{2}) > 0$, and the script computes the trace form's signature $(4,2)$ exactly by congruence diagonalisation (control: $\so(4)$ gives $(0,6)$). Suppose $\mathfrak{k} \subset \so(2,2)$ with $\mathfrak{k} \cong \su(2)$. The restriction of the trace form to $\mathfrak{k}$ is invariant, so it is a multiple of the Killing form of $\mathfrak{k}$. Take $X \in \mathfrak{k}$ nonzero; in $\mathfrak{k}\otimes\C \cong \mathfrak{sl}(2,\C)$, $X = \tfrac{I_{4}\alpha}{2}H$ for an $\mathfrak{sl}_{2}$-triple $(H, E, F)$ and real $\alpha \neq 0$, and on the faithful four-dimensional representation $H$ is diagonalisable with integer eigenvalues $m_{j}$, not all zero (imported: Fulton and Harris, Lecture 11, cited from memory). So $\tr(X^{2}) = -\tfrac{\alpha^{2}}{4}\sum_{j} m_{j}^{2} < 0$, and the trace form is negative definite on the three-dimensional $\mathfrak{k}$, which is impossible since its negative index on $\so(2,2)$ is $2$. A group of isometries of $(\mathcal{R}, Q)$ that is locally $SU(2)$ would have such a Lie algebra; and the holonomy algebra of a metric connection lies in $\so(\mathcal{R})$, whatever its curvature, so it contains no $\su(2)$.

(3) Let $K$ be symmetric and invariant under the $E_{ab}$. Invariance under $N_{R}$ gives $2K(v_{a}, v_{b}) = 0$ and $-2K(w_{a}, w_{b}) = 0$, so only the block $K(v_{a}, w_{b})$ survives; invariance under $E_{ab}$, $a \neq b$, makes that block a multiple of the identity. So $K = kQ$; the script solves the linear system exactly for $\gl(2,\R)$, for its $\mathfrak{sl}(2,\R)$ part and for $\so(2,2)$, finding a one-dimensional solution space each time. $Q$ on $\mathcal{R}$ has signature $(2,2)$. The form $\ip{\theta x}{y}$ is $\tfrac12$ times the identity in the basis $(v_{a}, w_{a})$ paired with its $\theta$-image, signature $(4,0)$ (script); it is not $\so(2,2)$-invariant, since a boost of a hyperbolic plane commutes with $\theta$ only when trivial.

(4) The lift of Proposition \claimref{standard-model.transported-frozen-point}(i) is $X(I) = n_{0}' + I + \tfrac{\Lambda}{6}n_{\infty}'$ for any null pair with $\ip{n_{0}'}{n_{\infty}'} = -1$; with $n_{0}' = -\lambda w_{4}$, $n_{\infty}' = 2\lambda^{-1}v_{4}$ the pairing is $-2\cdot\tfrac12 = -1$, and the rung component is $c_{\lambda} = -\lambda w_{4} + \tfrac{\Lambda}{3}\lambda^{-1}v_{4}$; at $\lambda = 1$ it is the proposition's $-w_{4} + \tfrac{\Lambda}{3}v_{4}$. Then $Q(c_{\lambda}) = 2(-\lambda)(\tfrac{\Lambda}{3}\lambda^{-1})\tfrac12 = -\tfrac{\Lambda}{3}$, and $\theta c_{\lambda} = -\lambda v_{4} + \tfrac{\Lambda}{3}\lambda^{-1}w_{4}$, so $\ip{\theta c_{\lambda}}{c_{\lambda}} = \tfrac12(\lambda^{2} + \tfrac{\Lambda^{2}}{9}\lambda^{-2})$, and subtracting $\tfrac{\Lambda}{3}$ leaves the square displayed. The script verifies both identities for twelve rational pairs $(\Lambda, \lambda)$. The boost $(n_{0}', n_{\infty}') \mapsto (e^{t}n_{0}', e^{-t}n_{\infty}')$ is an isometry of $\mathcal{R}$ fixing every vector of the ground space, so it changes no tractor norm; it changes $\theta$, and with it the $*$-norm.

(5) Let $\sigma$ be a conjugate-linear involution of $\mathcal{R}\otimes\C$ commuting with $U(2)$, with $Q$ real on its fixed set. The centre $\exp(I_{4}\varphi N_{R})$ multiplies $R_{2} = \mathrm{span}(v_{4}, v_{5})$ by $e^{I_{4}\varphi}$ and $\overline{R}_{2} = \mathrm{span}(w_{4}, w_{5})$ by $e^{-I_{4}\varphi}$, and $\sigma$ commutes with it, so $\sigma$ maps $R_{2}$ to $\overline{R}_{2}$; composing with complex conjugation of coefficients gives a $U(2)$-equivariant linear map of $R_{2}$, a scalar $\mu$ by Schur's lemma. So $\sigma(a\cdot v) = \mu\bar{a}\cdot w$, the fixed set is $\mathcal{R}_{\mu}$, and $Q(a\cdot v + \mu\bar{a}\cdot w) = \mu|a|^{2}$ is real exactly when $\mu$ is real. That $U(2)$ preserves each $\mathcal{R}_{\mu}$ is $\bar{X}(\mu\bar{a}) = \mu\overline{Xa}$; the script checks it for each generator and three values of $\mu$. The tower's $*$ sends $a\cdot v$ to $\bar{a}\cdot w$, so its fixed set is $\mathcal{R}_{1}$. Now $c_{\lambda} = a\cdot v + b\cdot w$ with $a = (\tfrac{\Lambda}{3}\lambda^{-1}, 0)$, $b = (-\lambda, 0)$, and $zc_{\lambda} \in \mathcal{R}_{\pm 1}$ requires $|zb| = |za|$, that is $\lambda^{2} = \Lambda/3$; the action of $U(2)$ preserves $|a|$ and $|b|$ separately, so no element of $U(2)\cdot\C^{\times}$ helps. Conversely at $\lambda^{2} = \Lambda/3$, $a = (\lambda, 0) = -b$, and $I_{4}c_{\lambda}$ has $I_{4}b = -I_{4}\lambda = \overline{I_{4}a}$, so it lies in $\mathcal{R}_{1}$. By (4), $\lambda^{2} = \Lambda/3$ is exactly where the $*$-norm reaches its minimum $|Q(c_{\lambda})|$. The script checks membership at a balanced pair and non-membership, through $|a| \neq |b|$, at unbalanced ones.

Reading, argued and not part of the statement. Standard-Model paper Proposition 9.29(2) (\claimref{standard-model.what-derived-what-not}) calls the coefficient one of $-w_{4}$ the rung unit and reads a weak scale of $246$ GeV as that unit being about $10^{44}$ Hubble units. By (4) the tractor geometry does not see this unit at all, and the positive norm a kinetic term would use does: with $H^{2} = \Lambda/3$ and $\lambda = H e^{s}$, $\ip{\theta c_{\lambda}}{c_{\lambda}} = H^{2}\cosh 2s$. A weak scale of the observed size needs $e^{2s}$ of order $(v/H)^{2} \approx 10^{88}$, while the one choice that makes the frozen point's rung component itself a value of a unitary doublet, $s = 0$ by (5), gives a $*$-norm of $H^{2}$, a weak scale equal to the Hubble rate. So the transported frozen point does not fix the weak scale through unitarity; the open number is the frame's boost, and by (2) and (3) the same frame is what a positive kinetic term and a compact weak connection need. Whether the frame is a datum of the series beside $N$ is a decision for the premise register, not settled here.
